\documentclass[UTF8]{ctexart}
\usepackage{amsmath, amssymb, amsthm}
\usepackage{geometry}
\usepackage{tikz}
\usepackage[colorlinks=true, linkcolor=blue, urlcolor=blue, citecolor=blue]{hyperref}
\geometry{margin=1in}

\newtheorem{definition}{定义}[section]
\newtheorem{proposition}[definition]{命题}
\newtheorem{remark}[definition]{注记}
\newtheorem{example}[definition]{例}
\newtheorem{fact}[definition]{事实}
\newtheorem{theorem}[definition]{定理}
\newtheorem{algorithm}[definition]{算法}

\title{《计算拓扑数据分析》第三章笔记：持续算法}
\author{“基于持续同调的拓扑特征识别及其应用”SRTP项目组}
\date{}

\begin{document}

\maketitle

\tableofcontents
\newpage

\section*{章节概览}
\addcontentsline{toc}{section}{章节概览}

本章聚焦于持续同调的\textbf{计算方面}。给定一个单纯滤过，我们的目标是计算其持续图。这等价于找出所有导致同调类出生与死亡的单纯形配对。我们将依次介绍：

\begin{enumerate}
    \item \textbf{瓶颈距离的计算}（\hyperref[sec:bottleneck]{BOTTLENECK 算法}）：通过二分图完美匹配计算两个持续图之间的距离；
    \item \textbf{配对算法}（\hyperref[sec:pairpersistence]{PAIRPERSISTENCE 算法}）：组合视角下的持续配对算法；
    \item \textbf{矩阵约简算法}（\hyperref[sec:matpersistence]{MATPERSISTENCE 算法}）：持续同调计算的核心算法，通过边界矩阵的列消元计算持续配对；
    \item \textbf{高效实现技巧}（\hyperref[sec:clearpersistence]{CLEARPERSISTENCE 算法}）：清除策略与上同调对偶优化；
    \item \textbf{零维持续同调算法}（\hyperref[sec:zeroperdg]{ZEROPERDG 算法}）：基于并查集的零维持续同调专用算法。
\end{enumerate}

\section{瓶颈距离的计算}
\label{sec:bottleneck}

在 3.2.1 节中我们定义了瓶颈距离，现在考虑其计算问题。

\subsection{问题转化}

给定两个持续图 \(\operatorname{Dgm}_p(\mathcal{F}_f)\) 和 \(\operatorname{Dgm}_p(\mathcal{F}_g)\)，设 \(A\) 和 \(B\) 分别为其非对角点集。为处理两个图点数不同的情况，我们将每个点投影到对角线上：

\begin{itemize}
    \item \(\bar{a}\)：点 \(a \in A\) 在对角线 \(\Delta\) 上的最近点；
    \item \(\bar{b}\)：点 \(b \in B\) 在对角线 \(\Delta\) 上的最近点；
    \item 构造扩展点集：\(\tilde{A} = A \cup \bar{B}\)，\(\tilde{B} = B \cup \bar{A}\)。
\end{itemize}

则瓶颈距离等于：
\[
d_b = \min_{\pi} \max_{a \in \tilde{A}} \|a - \pi(a)\|_\infty,
\]
其中 \(\pi\) 遍历 \(\tilde{A}\) 到 \(\tilde{B}\) 的所有双射。

\subsection{算法思路}

\begin{algorithm}[BOTTLENECK]
\label{alg:bottleneck}
\begin{enumerate}
    \item 计算所有可能的 \(O(n^2)\) 个点对 \(L_\infty\) 距离，排序得 \(\delta_0 \le \delta_1 \le \cdots \le \delta_{n'}\)；
    \item 二分搜索最小的 \(\delta\)，使得存在一个完美匹配；
    \item 对给定的 \(\delta\)，构造二分图 \(G = (\tilde{A} \cup \tilde{B}, E)\)，其中 \((a,b) \in E\) 当且仅当 \(\|a - b\|_\infty \le \delta\)；
    \item 检查是否存在完美匹配（Hopcroft-Karp 算法 \(O(n^{2.5})\)，利用几何结构可优化至 \(O(n^{1.5})\)）；
    \item 若存在，减小 \(\delta\)；否则增大 \(\delta\)。
\end{enumerate}
\end{algorithm}

\begin{remark}
\hyperref[alg:bottleneck]{BOTTLENECK 算法}的复杂度为 \(O(n^{1.5}\log n)\)，其中 \(n\) 是持续图中非对角点的数量。为避免排序所有 \(O(n^2)\) 个距离，可利用点的几何嵌入直接进行二分搜索。
\end{remark}

\section{持续配对算法（PAIRPERSISTENCE）}
\label{sec:pairpersistence}

\subsection{基本观察}

我们假设输入是\textbf{逐单形滤过}：
\[
\emptyset = K_0 \hookrightarrow K_1 \hookrightarrow \cdots \hookrightarrow K_n = K,
\]
其中 \(K_j \setminus K_{j-1} = \sigma_j\) 为单个单纯形。

\begin{fact}[添加单纯形时的两种可能]
\label{fact:positive-negative}
添加一个 \(p\)-单纯形 \(\sigma_j\) 时，恰好发生以下两种情况之一：
\begin{enumerate}
    \item \textbf{正单纯形（Positive）}：创建一个非边界的 \(p\)-闭链，\(\beta_p\) 增加 1；
    \item \textbf{负单纯形（Negative）}：销毁一个已有的 \((p-1)\)-闭链，\(\beta_{p-1}\) 减少 1。
\end{enumerate}
\end{fact}

\begin{remark}
正单纯形也称为“创建者（creator）”，负单纯形也称为“销毁者（destructor）”。每个负单纯形通过“最年轻优先”规则唯一配对一个正单纯形。
\end{remark}

\subsection{配对规则}

当一个 \(p\)-单纯形 \(\sigma_j\) 被添加时，检查其边界 \(\partial\sigma_j\) 是否是一个边界闭链（即已在当前同调群中为零）：

\begin{itemize}
    \item 若 \(\partial\sigma_j\) 已是边界，则 \(\sigma_j\) 创建了一个新的 \(p\)-维同调类（正单纯形）；
    \item 若 \(\partial\sigma_j\) 不是边界，则 \(\sigma_j\) 销毁了由某个更早的 \((p-1)\)-单纯形创建的类（负单纯形）。
\end{itemize}

为找到被销毁类的出生点，我们从 \(\partial\sigma_j\) 中取出\textbf{最年轻的}未配对的 \((p-1)\)-单纯形——这就是“最年轻优先”原则。

\subsection{算法描述}

\begin{algorithm}[PAIRPERSISTENCE]
\label{alg:pairpersistence}
\textbf{输入}：逐单形滤过的单纯形序列 \(\sigma_1, \sigma_2, \dots, \sigma_n\)

\textbf{输出}：每个单纯形的正/负标记及持续配对

\begin{enumerate}
    \item 对 \(j = 1\) 到 \(n\)：
    \begin{enumerate}
        \item 令 \(c := \partial_p \sigma_j\)；
        \item 令 \(\sigma_i\) 为 \(c\) 中最年轻的未配对的 \((p-1)\)-单纯形；
        \item 当 \(\sigma_i\) 已配对且 \(c\) 非空时：
        \begin{enumerate}
            \item 令 \(c'\) 为与 \(\sigma_i\) 配对的单纯形所销毁的闭链；
            \item \(c := c' + c\)（模 2 加法）；
            \item 更新 \(\sigma_i\) 为 \(c\) 中最年轻的未配对的 \((p-1)\)-单纯形；
        \end{enumerate}
        \item 若 \(c\) 非空：\(\sigma_j\) 为负单纯形，生成配对 \((\sigma_i, \sigma_j)\)；
        \item 否则：\(\sigma_j\) 为正单纯形。
    \end{enumerate}
\end{enumerate}
\end{algorithm}

\section{矩阵约简算法（MATPERSISTENCE）}
\label{sec:matpersistence}

\subsection{过滤边界矩阵}

\begin{definition}[过滤边界矩阵]
\label{def:filtered-boundary}
设滤过 \(\emptyset = K_0 \hookrightarrow K_1 \hookrightarrow \cdots \hookrightarrow K_n = K\)，单纯形按滤过顺序排列为 \(\sigma_1, \dots, \sigma_n\)。\textbf{过滤边界矩阵} \(D\) 是 \(n \times n\) 矩阵，其中
\[
D[i,j] = 
\begin{cases}
1, & \sigma_i \in \partial \sigma_j,\\
0, & \text{否则}.
\end{cases}
\]
即第 \(j\) 列表示 \(\sigma_j\) 的边界（在 \(\mathbb{Z}_2\) 系数下）。
\end{definition}

\begin{remark}
因 \(\partial^2 = 0\)，边界矩阵满足 \(D^2 = 0\)（在适当的意义下）。
\end{remark}

\subsection{约简矩阵与低行索引}

\begin{definition}[低行索引]
\label{def:low}
对矩阵 \(A\)，\(\operatorname{low}_A[j]\) 表示第 \(j\) 列中\textbf{最后一个} \(1\) 所在的行索引（即最大的非零行索引）。若列为空，则 \(\operatorname{low}_A[j] = -1\)。
\end{definition}

\begin{definition}[约简矩阵]
\label{def:reduced}
矩阵 \(A\) 称为\textbf{约简的（reduced）}，若对任意 \(j \neq j'\)，\(\operatorname{low}_A[j] \neq \operatorname{low}_A[j']\)，即所有非空列的低行索引互不相同。
\end{definition}

\begin{fact}
约简矩阵的非零列在 \(\mathbb{Z}_2\) 上线性无关。
\end{fact}

\subsection{核心算法}

\hyperref[alg:matpersistence]{MATPERSISTENCE 算法}通过\textbf{从左到右的列加法}将过滤边界矩阵约简为约简形式。列加法 \( \operatorname{col}_j \leftarrow \operatorname{col}_j + \operatorname{col}_i\)（\(i < j\)）对应于将第 \(i\) 列加到第 \(j\) 列。

\begin{algorithm}[MATPERSISTENCE]
\label{alg:matpersistence}
\textbf{输入}：过滤边界矩阵 \(D\)（列和行按滤过顺序排列）

\textbf{输出}：约简矩阵，其中每列要么为空，要么有唯一的 \(\operatorname{low}\) 值

\begin{enumerate}
    \item 对 \(j = 1\) 到 \(|D|\)：
    \begin{enumerate}
        \item 当存在 \(j' < j\) 使得 \(\operatorname{low}[j'] = \operatorname{low}[j]\) 且 \(\operatorname{low}[j] \neq -1\)：
        \begin{enumerate}
            \item \(\operatorname{col}[j] \leftarrow \operatorname{col}[j] + \operatorname{col}[j']\)
        \end{enumerate}
        \item 若 \(\operatorname{low}[j] \neq -1\)，则生成配对 \((\sigma_{\operatorname{low}[j]}, \sigma_j)\)
    \end{enumerate}
\end{enumerate}
\end{algorithm}

\begin{remark}
算法的核心逻辑是：当两列的低行索引冲突时，将左侧列加到右侧列上，直到冲突消除。这本质上是在矩阵上执行“最年轻优先”配对。
\end{remark}

\subsection{配对定理}

\begin{theorem}[持续配对判定定理]
\label{thm:persistence-pairing}
设 \(D\) 为过滤边界矩阵，\(R = DV\) 为 \(D\) 的约简形式（\(V\) 为上三角矩阵）。则
\[
\sigma_i \text{ 与 } \sigma_j \text{ 形成持续配对} \iff \operatorname{low}_R[j] = i.
\]
\end{theorem}

\begin{remark}
该定理表明，持续配对完全由约简矩阵中的低行索引决定，且与具体的约简过程无关。
\end{remark}

\subsection{复杂度分析}

\begin{itemize}
    \item \hyperref[alg:matpersistence]{MATPERSISTENCE 算法}：\(O(n^3)\)，其中 \(n\) 为单纯形总数；
    \item 利用稀疏矩阵表示可优化；
    \item 使用快速矩阵乘法可达到 \(O(n^\omega)\)，其中 \(\omega \in [2, 2.373)\) 为矩阵乘法指数。
\end{itemize}

\section{高效实现技巧}
\label{sec:clearpersistence}

\subsection{清除策略（Clearing）}

\begin{remark}[核心观察]
若一个 \(p\)-单纯形 \(\sigma_j\) 是正单纯形，它将在更高维度的处理中与某个 \((p+1)\)-单纯形配对。因此在处理 \(p\) 维边界矩阵时，可\textbf{跳过这些已知的正单纯形}，不必将其约简为零列。
\end{remark}

\begin{algorithm}[CLEARPERSISTENCE]
\label{alg:clearpersistence}
\textbf{输入}：按维度排列的边界矩阵 \(D_1, D_2, \dots, D_d\)

\textbf{输出}：约简矩阵，每个负单纯形的列具有唯一的低行索引

\begin{enumerate}
    \item \hyperref[alg:matpersistence]{MATPERSISTENCE}(\(D_d\))；
    \item 对 \(i = d-1\) 到 \(1\)：
    \begin{enumerate}
        \item 对 \(D_i\) 的每一列 \(j\)：
        \begin{enumerate}
            \item 若 \(\sigma_j\) 在处理 \(D_{i+1}\) 时已配对，则跳过该列；
            \item 否则，执行与 \hyperref[alg:matpersistence]{MATPERSISTENCE} 相同的列约简；
            \item 若约简后非空，生成配对。
        \end{enumerate}
    \end{enumerate}
\end{enumerate}
\end{algorithm}

\begin{remark}
清除策略在实际计算中（特别是 Rips 滤过）能大幅提升效率，因为高维正单纯形的数量通常远大于负单纯形。
\end{remark}

\subsection{扭曲矩阵（上同调算法）}

\begin{definition}[扭曲矩阵]
\label{def:twisted}
矩阵 \(D_p\) 的\textbf{扭曲矩阵} \(D_p^*\) 定义为 \(D_p\) 的“反转置”（先转置，再反转行列顺序）：
\[
D_p^*(i, j) = D_p(n + 1 - j, m + 1 - i).
\]
\end{definition}

\begin{proposition}[扭曲矩阵的对偶性]
\label{prop:twisted-duality}
在 \(D_p\) 中配对 \((\sigma^{p-1}, \sigma^p)\) 当且仅当在 \(D_p^*\) 中配对 \((\sigma^p, \sigma^{p-1})\)。
\end{proposition}

\begin{remark}
这一对偶性允许我们选择计算效率更高的方向：若高维单纯形较多，使用扭曲矩阵从低维往高维处理可能更优。
\end{remark}

\section*{算法总结}
\addcontentsline{toc}{section}{算法总结}

\begin{tabular}{|c|c|c|c|}
\hline
\textbf{算法} & \textbf{输入} & \textbf{输出} & \textbf{核心思想} \\
\hline
\hyperref[alg:bottleneck]{BOTTLENECK} & 两个持续图 & 瓶颈距离 & 二分图完美匹配 + 二分搜索 \\
\hline
\hyperref[alg:pairpersistence]{PAIRPERSISTENCE} & 单纯形序列 & 持续配对 & 最年轻优先原则 \\
\hline
\hyperref[alg:matpersistence]{MATPERSISTENCE} & 过滤边界矩阵 & 约简矩阵 + 配对 & 从左到右列消元 \\
\hline
\hyperref[alg:clearpersistence]{CLEARPERSISTENCE} & 各维边界矩阵 & 约简矩阵 & 跳过已知正单纯形 \\
\hline
\hyperref[alg:zeroperdg]{ZEROPERDG} & 1-复形 + 顶点函数 & 零维持续配对 & 并查集维护连通分量 \\
\hline
\end{tabular}

\section{零维持续同调算法（ZEROPERDG）}
\label{sec:zeroperdg}

对于一般持续同调算法，已知最优运行时间为 \(O(n^\omega)\)，其中 \(n\) 是滤过中单纯形的总数。然而，对于 \textbf{零维持续同调}（即持续图 \(\operatorname{Dgm}_0 f\)），存在更为高效的专用算法。

\subsection{核心观察}

计算 \(\operatorname{Dgm}_0 f\) 仅需复形 \(K\) 的 \textbf{1-骨架}（即顶点集 \(V\) 和边集 \(E\)），因为零维同调仅关心连通分量的变化。

设顶点按函数值非递减排序：
\[
f(v_1) \le f(v_2) \le \cdots \le f(v_n).
\]
令 \(K_i\) 为前 \(i\) 个顶点的下星之并，即 \(K_i = \bigcup_{j \le i} \overline{\operatorname{Lst}}(v_j)\)。零维同调 \(H_0(K_i)\) 的本质信息就是 \(K_i\) 中连通分量的数量与组成。

\subsection{三种情况}

当处理顶点 \(v_j\) 时，考察其下星 \(\operatorname{Lst}(v_j)\)：

\begin{enumerate}
    \item \textbf{情况 1}：\(\operatorname{Lst}(v_j) = \{v_j\}\)。此时 \(v_j\) 在 \(K_j\) 中创建一个新的连通分量，\(v_j\) 为\textbf{创建者}。
    
    \item \textbf{情况 2}：\(\operatorname{Lst}(v_j)\) 中的所有边都连接到 \(K_{j-1}\) 中的\textbf{同一个连通分量} \(C\)。此时 \(C\) 扩展包含 \(v_j\) 及其关联边，但 \(H_0(K_{j-1}) \cong H_0(K_j)\)，无拓扑变化。
    
    \item \textbf{情况 3}：\(\operatorname{Lst}(v_j)\) 中的边连接到 \(K_{j-1}\) 中的\textbf{多个连通分量} \(C_1, C_2, \dots, C_r\)（\(r \ge 2\)）。此时这些分量合并为一个新分量：
    \[
    C' = C_1 \cup C_2 \cup \cdots \cup C_r \cup \operatorname{Lst}(v_j).
    \]
    包含映射 \(K_{j-1} \hookrightarrow K_j\) 诱导一个满射同态 \(\xi: H_0(K_{j-1}) \to H_0(K_j)\)，且
    \[
    \beta_0(K_j) = \beta_0(K_{j-1}) - (r - 1),
    \]
    即 \(r-1\) 个分量被销毁，仅一个存活为 \(C'\)。
\end{enumerate}

\subsection{配对规则}

\begin{proposition}[Proposition 3.17]
\label{prop:zerod-merge}
设情况 3 发生，\(\operatorname{Lst}(v_j)\) 中的边合并了 \(K_{j-1}\) 中的分量 \(C_1, \dots, C_r\)。令 \(v_{k_i}\) 为分量 \(C_i\) 中的\textbf{全局最小值顶点}（即 \(f(v_{k_i})\) 最小），并设
\[
f(v_{k_1}) \le f(v_{k_2}) \le \cdots \le f(v_{k_r}).
\]
则顶点 \(v_j\) 恰好参与 \(r-1\) 个零维持续配对：
\[
(v_{k_2}, v_j),\; (v_{k_3}, v_j),\; \dots,\; (v_{k_r}, v_j),
\]
对应持续图中的点
\[
(f(v_{k_2}), f(v_j)),\; (f(v_{k_3}), f(v_j)),\; \dots,\; (f(v_{k_r}), f(v_j)).
\]
\end{proposition}

\begin{remark}[直观理解]
当多个分量合并时，考虑 0-闭链：
\[
c_2 = v_{k_2} + v_{k_1},\quad c_3 = v_{k_3} + v_{k_1},\quad \dots,\quad c_r = v_{k_r} + v_{k_1}.
\]
它们在 \(H_0(K_{j-1})\) 中独立，各自在 \(K_{k_i}\) 处创建，但在进入 \(K_j\) 时全部变为平凡（被销毁）。只有最“古老”的分量 \(v_{k_1}\) 对应的类 \(c_1 = v_{k_1} + c\) 继续存活，不参与死亡配对。
\end{remark}

\subsection{算法描述}

基于 \hyperref[prop:zerod-merge]{Proposition 3.17}，零维持续同调的计算无需矩阵约简，只需维护各 \(K_i\) 的连通分量信息，支持：
\begin{itemize}
    \item 查询顶点所属分量的代表元；
    \item 合并多个分量的操作。
\end{itemize}

这正是 \textbf{并查集（Union-Find）} 数据结构的典型应用。

\begin{algorithm}[ZEROPERDG]
\label{alg:zeroperdg}
\textbf{输入}：1-复形 \(K = (V, E)\)，顶点函数 \(f: V \to \mathbb{R}\)

\textbf{输出}：生成 \(\operatorname{Dgm}_0(f)\) 的顶点配对

\begin{enumerate}
    \item 按 \(f\) 值升序排序顶点：\(v_1, \dots, v_n\)
    \item 对 \(j = 1\) 到 \(n\)：
        \begin{enumerate}
            \item \(\operatorname{CREATESET}(v_j)\)
            \item 令 \(flag \leftarrow 0\)
            \item 对 \(\operatorname{Lst}(v_j)\) 中每条边 \((v_k, v_j)\)：
                \begin{itemize}
                    \item 若 \(flag = 0\)，则执行 \(\operatorname{UNION}(v_k, v_j)\) 并置 \(flag \leftarrow 1\)；
                    \item 否则，若 \(\operatorname{FINDSET}(v_k) \ne \operatorname{FINDSET}(v_j)\)，则：
                        \begin{itemize}
                            \item 令 \(\ell_1 \leftarrow \operatorname{REPSET}(v_k)\)，\(\ell_2 \leftarrow \operatorname{REPSET}(v_j)\)
                            \item \(\operatorname{UNION}(v_k, v_j)\)
                            \item 输出配对 \((\arg\max\{f(\ell_1), f(\ell_2)\}, v_j)\)
                        \end{itemize}
                \end{itemize}
        \end{enumerate}
    \item 对每个最终的分量 \(C\)，输出配对 \((\operatorname{REPSET}(C), \infty)\)
\end{enumerate}
\end{algorithm}

\begin{remark}
并查集操作说明：
\begin{itemize}
    \item \(\operatorname{CREATESET}(x)\)：创建仅含 \(x\) 的新集合
    \item \(\operatorname{FINDSET}(x)\)：返回 \(x\) 所在集合的代表元
    \item \(\operatorname{UNION}(x, y)\)：合并 \(x\) 与 \(y\) 所在的集合
    \item \(\operatorname{REPSET}(x)\)：返回 \(x\) 所在集合中 \(f\) 值最小的顶点（即代表元）
\end{itemize}
\end{remark}

\subsection{复杂度分析}

\begin{itemize}
    \item 排序顶点：\(O(n \log n)\)
    \item 并查集操作（CREATE、FIND、UNION、REPSET）次数：\(O(n + m)\)
    \item 总时间复杂度：\(O(n \log n + m \alpha(n))\)，其中 \(\alpha(n)\) 为反阿克曼函数（增长极慢，可视为常数）
\end{itemize}

\begin{theorem}[Theorem 3.18]
对于 PL-函数 \(f: |K| \to \mathbb{R}\)，零维持续图 \(\operatorname{Dgm}_0 f\) 可由 \hyperref[alg:zeroperdg]{ZEROPERDG 算法} 在 \(O(n \log n + m \alpha(n))\) 时间内计算，其中 \(n\) 和 \(m\) 分别为 \(K\) 的顶点数和边数。
\end{theorem}

\subsection{与最小生成树的联系}

若将 \(K\) 的 1-骨架视为加权图 \(G = (V, E)\)，边权重设为
\[
w(u, v) = \max\{f(u), f(v)\},
\]
则 \hyperref[alg:zeroperdg]{ZEROPERDG 算法} 本质上等价于在 \(G\) 上运行 \textbf{Kruskal 最小生成树算法}。

当一条边 \(e = (u, v)\) 连接两个不同分量时，找出两个分量的最小顶点 \(\ell_1, \ell_2\)，将 \(e\) 与其中 \(f\) 值较大的那个顶点配对。然后将所有顶点-边配对 \((u, e)\) 转换为顶点-顶点配对 \((u, v)\)，其中 \(e \in \operatorname{Lst}(v)\)；形如 \((u, u)\) 的本地配对被丢弃。

\subsection{图滤过情形}

\hyperref[alg:zeroperdg]{ZEROPERDG 算法} 可直接推广到图滤过：按滤过顺序处理顶点和边，用并查集维护连通分量。每条连接两个不同分量的边 \(e = (u, v)\)，与两个分量中较年轻的顶点配对。未配对的顶点和边分别提供零维和一维持续图中的无限条。

\begin{remark}
在图滤过情形下，顶点排序由输入滤过隐式给出，因此复杂度可降至 \(O(n \alpha(n))\)。
\end{remark}

\end{document}